% Ч1.1 — Дроби: основное свойство, перевод десятичных дробей в обыкновенные
% Старая методичка: стр. 216 (старое задание 19)

\begin{tema}{Ч1.1}
\opr{Основное свойство дроби.} Если числитель и знаменатель умножить или разделить
на одно и то же число (выражение), не равное нулю, получится равная дробь:
\[
  \frac{a}{b}=\frac{a\cdot k}{b\cdot k}, \qquad b\ne 0,\; k\ne 0.
\]
На этом держится сокращение дробей и приведение к общему знаменателю — и для чисел,
и для алгебраических выражений.

\opr{Десятичная дробь → обыкновенная.} Пишем число без запятой в числителе,
а в знаменателе — 10, 100, 1000…, и сокращаем:
\[
  0{,}125=\frac{125}{1000}=\frac18, \qquad 1{,}125=1+\frac18=\frac98.
\]

\begin{pfTablica}{*{7}{c}}
$0{,}5$ & $0{,}25$ & $0{,}75$ & $0{,}2$ & $0{,}125$ & $0{,}375$ & $0{,}625$\\
$\frac12$ & $\frac14$ & $\frac34$ & $\frac15$ & $\frac18$ & $\frac38$ & $\frac58$\\
\end{pfTablica}

\begin{primer}
$1{,}125^4=\left(\frac98\right)^4=\frac{6561}{4096}$ — в обыкновенных дробях
считать степени проще, чем в десятичных.
\end{primer}

\begin{oshibka}
Сокращают слагаемые: $\frac{a+3}{3}=a$ — неверно.
Сокращать можно только \emph{множители}: $\frac{3a}{3}=a$.
\end{oshibka}
\end{tema}
